\section{Problem 3.}

A roll of paper towels forms a hollow cylinder with outer diameter 5”, inner diameter 1 3/4”, and height 11”.  
If the roll contains 120 sheets 11” by 10.4” of uniform thickness with no space between adjacent sheets, 
find the thickness of a sheet to the nearest ten-thousandth of an inch.

\begin{solution}
This reminds me of adding up onion rings of circumference to 
obtain the area of a circle. The approach is to either think of it as 
volumes and use the thickness surface area formula or to remove a dimension 
and think of wrapping string around a circle with subdivisions relating to 
the thickness. Let us go with the latter.

We have a height component that is redundant. This will simplify it to basically
a string of length 10.4” wrapped around a hollow tube until we get to the outer diameter
of 5”. Let us now create a formula that can calculate the total unwrapped circumference
based on summations. 

First let us define the intervals.
\begin{align*}
    & r_{\text{range}} = \frac{5}{2} - \frac{1.75}{2} = \frac{13}{8} \\
    & \text{interval}  = \frac{r_{\text{range}}}{\text{subdivisions}}
\end{align*}

Now let us create our function using summations.
\begin{equation*}
    C(x) = 2\pi r_{in}+\sum_{n=1}^{x}2\pi\left(r_{in}+n\frac{r_{\text{range}}}{x+1}\right), \ x \in \mathbb{N}
\end{equation*}

With this, we just need to solve for when the function $C(x)$ is equal to 
the total wrapping of 120 sheets of 10.4” length string.

\begin{align*}
    & 120 \cdot 10.4 = 1248 \\
 \Longrightarrow   & C(x) = 1248
\end{align*}

Graphing this out results in $x \approx 116.61291$. Since subdivisions can only be 
natural numbers, we round up which makes $x = 117$. This gets us an approximate thickness of
$\frac{r_{\text{range}}}{117} = \boxed{0.013\overline{8}”}$.

Comparing this to the volumne thickness method, 
I want to test and compare the wrapping answer the thickness volume method. 
To start, we must calculate the volume from annulus area multiplied by the height.
Then, we will divide that by the surface area of the the 120 sheets to get our thickness.

\begin{align*}
    V_{\text{roll}} &= \left(\pi r_{\text{out}}^2 - \pi r_{\text{in}}^2\right) \cdot h \\
    V_{\text{sheets}} &= 120(h \cdot 10.4) \cdot t \\
    V_{\text{sheets}} &= V_{\text{roll}} \\
    120(h \cdot 10.4) \cdot t &= \left(\pi r_{\text{out}}^2 - \pi r_{\text{in}}^2\right) \cdot h  \\
    t &= \frac{\left(\pi r_{\text{out}}^2 - \pi r_{\text{in}}^2\right) \cdot \cancel{h}}{120(\cancel{h} \cdot 10.4)} \\
      &= \frac{\pi\cdot\left(2.5\right)^{2}-\pi\left(0.875\right)^{2}}{120\cdot10.4} \\
    t &\approx \boxed{0.1381"}
\end{align*}

So 2 decimal places of accuracy which matches our significant figures. Not too shabby.
Would say that the initial method was a bit overdone compared the volume approach.

\end{solution}
